\newpage

\thispagestyle{empty}

\begin{center}
\textbf{I am deeply indebted to the following people, who all contributed in some way to my success in -- and/or well-being during -- the writing
of this thesis:}\bigskip

First and foremost my parents, for their unconditional support in all my endeavours, no matter how seemingly misguided.\bigskip

My wonderful girlfriend Tina, for her unrestrained love and affection, and for happily playing \emph{Tropico} next to me when I need to get work done while being with her.\bigskip

My band mates in \emph{Vroudenspil}, for providing me with a desperately needed outlet outside of my profession, besides being genuinely lovely friends to have. \emph{\# ShamelessPlug}\bigskip

Many researchers I had the pleasure to meet, and who $\langle$\emph{ any $x\in\mathcal P(\{$inspired, taught, advised, humoured, supported, entertained, just generally have been extremely nice to$\})\backslash\{\varnothing\}$}$\;\rangle$ me in various ways; including but not limited to: Jacques Carette, Merlin Carl, Claudio Sacerdoti Coen, Paul-Olivier Dehaye, Bill Farmer, Michael Junk, Cezary Kaliszyk, Peter Koepke, Andrea Kohlhase, Thomas Koprucki, Markus Pfeiffer, Natarajan Shankar, Nicolas Thiéry, Makarius Wenzel, ...\bigskip

My colleagues at KWARC, for contributing to establishing a productive, and yet unusually comfortable and fun work environment: Katja Ber\v ci\v c, Jonas Betzendahl, Deyan Ginev, Mihnea Iancu, Constantin Jucovschi, Theresa Pollinger, Max Rapp, Frederik Schaefer, Tom Wiesing. 

In particular Florian Rabe, who acted as a second advisor, especially in technical matters, and helped me literally every step of the way.\bigskip

And last but not least Michael Kohlhase, for being an exceptionally passionate, caring and inspiring advisor, and for avoiding without exception all negative tropes about PhD advisors that make up the central plots of countless horror stories that traumatized post grads tell each other huddled around bonfires made from $n$th revisions of (ultimately failed) thesis drafts.
\end{center}\bigskip

\newpage

\thispagestyle{empty}

\begin{quotation} \centering\small%\tiny
%``Be warned, there will be lots of similes and obscure aphorisms, which start well but end up making no sense.

%So listen up, or you’ll get lost, like a blind man clapping in a pharmacy trying to echo-locate the contact lens fluid.''\medskip

%``Searching for meaning [in life] is like searching for a rhyme scheme in a cookbook: You won’t find it, and it will bugger up your soufflé.''\medskip

\textbf{``You don’t have to have a dream.}
Americans on talent shows always talk about their dreams. Fine, if you have something you’ve always wanted to do, dreamed of, like, in your heart, go for it. After all, it’s something to do with your time, chasing a dream. And if it’s a big enough one, it’ll take you most of your life to achieve, so by the time you get to it and are staring into the abyss of the meaninglessness of your achievement you’ll be almost dead, so it won’t matter.

I never really had one of these dreams, and so I advocate passionate dedication to the pursuit of short-term goals. Be micro-ambitious. Put your head down and work with pride on whatever is in front of you. You never know where you might end up. Just be aware, the next worthy pursuit will probably appear in your periphery, which is why you should be careful of long-term dreams. If you focus too far in front of you, you won’t see the shiny thing out the corner of your eye.''\medskip

\textbf{``Don’t seek happiness.} 	
Happiness is like an orgasm: If you think about it too much it goes away. Keep busy and aim to make someone else happy, and you might find you get some as a side effect.''\medskip

\textbf{ ``Be hard on your opinions.}
%
A famous bon mot asserts that opinions are like arseholes in that everyone has one. There is great wisdom in this, but I would add that opinions differ significantly from arseholes in that yours should be constantly and thoroughly examined.''\bigskip
	
\textbf{``Remember it’s all luck.}	
You are lucky to be here. You are incalculably lucky to be born, and incredibly lucky to be brought up by a nice family that helped you get educated and encouraged you to go to uni.

Or, if you were born into a horrible family, that's unlucky and you have my sympathy, but you are still lucky. Lucky that you happen to be made of the sort of DNA that went on to make the sort of brain which, when placed in a horrible childhood environment, would make decisions that meant you ended up eventually graduating uni. Well done, you, for dragging yourself up by your shoelaces. But you were lucky. You didn't create the bit of you that dragged you up. They're not even your shoelaces.

I suppose I worked hard to achieve whatever dubious achievements I've achieved, but I didn't make the bit of me that works hard, any more than I made the bit of me that ate too many burgers instead of attending lectures when I was [at uni].

Understanding that you can't \emph{truly} take credit for your successes, nor \emph{truly} blame others for their failures, will humble you and make you more compassionate. 

Empathy is intuitive, but it is also something you can work on intellectually.''\bigskip
%	
%%\medskip
%	
%%``Please don’t make the mistake of thinking the arts and sciences are at odds with one another. That is a recent, stupid and damaging idea.
%
%%You don’t have to be unscientific to make beautiful art or to write beautiful things. [...] You don’t need to be superstitious to be a poet. You don’t need to hate GM technology to care about the beauty of the planet. You don’t have to claim a soul to promote compassion.''\medskip
%

%``Arts degrees are awesome, and they help you find meaning where there is none.
%And let me assure you: \textbf{There is none.}
%Don't go looking for it. 
%Searching for meaning is like searching for a rhyme scheme in a cookbook: You won't find it, and it will bugger up your soufflé.[...]
%[This is] not a flippant assertion. I think it’s absurd, the idea of seeking meaning in the set of circumstances that happens to exist after 13.8 billion years worth of unguided events. Leave it to humans to think the universe has a purpose for them.
%
%However, I’m no nihilist. I’m not even a cynic. I am actually rather romantic. And here’s my idea of romance: You will soon be dead.
%
%Life will sometimes seem long, and tough, and God, it’s tiring. And you will sometimes be happy, and sometimes sad, and then you’ll be old, and then you’ll be dead.
%
%There is only one sensible thing to do with this empty existence, and that is: \textbf{fill it.}
%[...]
%
%And in my opinion, until I change it, life is best filled by learning as much as you can, about as much as you can, taking pride in whatever you’re doing, having compassion, sharing ideas, running, being enthusiastic. And then there’s love and travel and wine and sex and art and kids and giving and mountain climbing, but you know all that stuff already.
%
%It’s an incredibly exciting thing, this one, meaningless life of yours.''\medskip
	
	\centering - Tim Minchin, Graduation Speech, University of Western Australia, 2013\footnote{\url{http://www.news.uwa.edu.au/201309176069/alumni/tim-minchin-stars-uwa-graduation-ceremony}}
\end{quotation}

\newpage

